%%-- Other slides just in case--
%% ── Extra slides (will delete after) ─────────────────────────────────────────
\begin{frame}{Modeling — What changed vs Logistic Regression}
\presenterphoto{sam_ml}
  \begin{block}{1. Tree models flip the feature-regime story}
    \footnotesize Logistic regression \emph{got worse} as features were added
    (1.06 $\to$ 1.14 $\to$ 1.22): multicollinearity and curse of
    dimensionality.  Tree models \emph{improve} with more features
    (RF: 1.26 $\to$ 1.06 $\to$ 1.05) — they handle correlated inputs natively.
  \end{block}

  \begin{block}{2. Player features add real signal}
    \footnotesize The full player regime (C) is the best for every tree model.
    Confirms the project hypothesis: player-level rolling stats add
    predictive power once the model can use them.
  \end{block}

  \begin{block}{3. Ensembling extracts a final gain}
    \footnotesize RF and XGB make different mistakes; averaging their
    calibrated probabilities beats either alone (1.04 vs 1.05/1.05 in regime C).
  \end{block}
\end{frame}


%% ── KNN: how we chose k (extra slides to answer questions) ────────────────────────────────────────
\begin{frame}{KNN — Choosing $k$}

  \footnotesize $k$ was not hand-picked. The full grid
  $k \in \{5, 10, 15, 20, 30, 50\} \times$ \{uniform, distance\}
  $\times$ \{Manhattan, Euclidean\} was searched under the same
  \texttt{TimeSeriesSplit(5)} protocol and the lowest-log-loss
  combination was kept.

  \vspace{0.4em}
  \begin{columns}[T]
    \column{0.5\textwidth}
    \begin{block}{Why those bounds?}
      \scriptsize
      \textbf{MIN $k=5$:} smaller $k$ gives degenerate probabilities
      (denominators of 2--3 only output $\{0, 0.33, 0.5, 0.66, 1\}$)
      — log-loss punishes that hard.\\[0.3em]
      \textbf{MAX $k=50$:} the smallest CV training fold is $\sim$75 rows;
      sklearn requires $k \le$ fold size. We initially tried $k=100$
      and got runtime warnings, so we capped at 50.
    \end{block}

    \column{0.48\textwidth}
    \begin{block}{Winners (per regime)}
      \scriptsize
      \setlength{\tabcolsep}{4pt}
      \begin{tabular}{lcc}
        Regime & $k$ & metric \\
        \hline
        A: team   & 50 & Euclidean ($p{=}2$) \\
        B: +FBref & 50 & Manhattan ($p{=}1$) \\
        C: +player& 50 & Manhattan ($p{=}1$) \\
      \end{tabular}\\[0.4em]
      All three regimes \textbf{independently} picked $k=50$ with
      distance-weighting. The Minkowski exponent drifted from
      $p{=}2$ to $p{=}1$ as features grew — consistent with the
      textbook result that L1 is more robust in higher dimensions.
    \end{block}
  \end{columns}

  \vspace{0.4em}
  \footnotesize\textbf{Why $k=50$ makes domain sense:} distance-weighting
  means the closest 5--10 historical matchups still drive the
  prediction; the further 40 just provide a regularising prior.
  Football is noisy, so averaging $\sim$12\% of the historical sample
  smooths randomness without washing out the matchup-specific signal.
\end{frame}